\makereporttitle{2025}{februarie}

\makereportpreamble

\vspace{0.5em}

În această lună am reușit să fac progrese semnificative pe proiectul de cercetare în care studiez apariția haosului cuantic în modelul cuadrupol bosonic al oscilațiilor nucleare. Extinzând codul de simulare numerică a dinamicii sistemului (în limita semiclasică), am reușit să identific noi comportamente interesante ale distribuției exponenților Lyapunov (pentru condiții inițiale distribuite uniform aleator în spațiul fazelor) atunci când (hiper)parametrul \(B\) variază în jurul valorii \(0.1\). Pentru energii suficient de înalte, se poate observa că exponentul Lyapunov maximal crește abrupt iar apoi urmează o curbă exponențială.

\begin{figure}[htbp]
    \centering
    \includegraphics[width=0.8\linewidth]{reports/2025/images/2/B-parameter-lyapunov-exponent-dependence}
    \caption*{Dependența exponentului Lyapunov maximal de valoarea parametrului \(B\)}
\end{figure}

În următoarea perioadă voi continua să analizez și să interpretez datele sintetice obținute, urmând apoi să fac și alte experimente cu modelul dezvoltat. Totodată, am reușit să extind și să optimizez și codul de diagonalizare exactă a Hamiltonianului cuantic. Voi începe să clasific stăriile proprii (vectorii proprii) ale Hamiltonianului în funcție de reprezentările ireductibile ale grupului de simetrie a sistemului, \(C_{3 v}\), la care corespund, pentru a putea apoi să reprezint grafic și să analizez distribuția lor și a spațiilor dintre ele. În acest sens, am început să parcurg cartea ``\emph{Representations and Characters of Groups}'' a lui Gordon James și Martin Liebeck.

În paralel, am continuat să particip la manifestări științifice și să învăț noțiuni importante pe cont propriu. De exemplu, am parcurs modulele de învățare și am rezolvat exercițiile lăsate ca temă din cursul online ``\emph{Mathematical and Computational Methods for Physics}'' oferit de Georgetown University pe platforma edX. Am continuat să lucrez exercițiile de geometrie diferențială și Riemanniană lăsate ca temă în cursul de ``\emph{Riemannian Geometry}'', pe care le-am completat cu parcurgerea unor materiale online despre teoria relativității generale. Am rezolvat în continuare exerciții de calcul numeric lăsate ca temă la cursul ``\emph{Metode numerice pentru probleme neliniare și optimizare}'' și m-am folosit de experiența dobândită din acestea pentru a progresa și cu proiectul de cercetare pe construcția iterativă a matricelor de precondiționare pentru sisteme liniare mari.

De asemenea, am ajutat la operaționalizarea noului \emph{Advanced Computing Center} a Universității din București, colaborând cu drd.\ Petre Boboc și cu dl.\ Marian Ancuța de la DITC-UB pentru a pune în funcțiune sistemul de calcul.

\textbf{Observațiile conducătorului științific}:

\vspace{3em}

\makereportpostamble{28 februarie 2025}
